\documentclass[11pt]{article}

\usepackage[a4paper,margin=2.2cm]{geometry}
\usepackage{amsmath,amssymb}
\usepackage{bm}
\usepackage{siunitx}
\usepackage{booktabs}
\usepackage{xcolor}
\usepackage[colorlinks=true,linkcolor=blue]{hyperref}

\newcommand{\code}[1]{\texttt{#1}}

\title{\textbf{Apollo Lunar Module --- One-Shot Motion-Planning OCP}\\[2pt]
\large Code Guide \;(\code{apollo\_full.py})}
\author{}
\date{}

\begin{document}
\maketitle
\vspace{-2.5em}

\section*{Overview}

\code{apollo\_full.py} is a single self-contained script that plans a soft lunar
landing. It builds \emph{one} open-loop optimal control problem (OCP), solves the
resulting nonlinear program (NLP) \emph{once}, appends a short analytic
engine-cut-off phase, and plots the result. The method is \textbf{direct multiple
shooting} with an \textbf{RK4} integrator, assembled with CasADi's \code{Opti}
modelling layer and solved by \textbf{IPOPT}. The end-to-end pipeline is

\[
\text{parameters}\;\longrightarrow\;
\boxed{\code{solve\_landing}}\;\longrightarrow\;
\boxed{\code{cutoff\_freefall}}\;\longrightarrow\;
\text{report + plots}.
\]

\section*{File layout}

The file follows six banner-marked sections:
\begin{enumerate}\itemsep1pt
  \item \textbf{Parameters} --- three \code{@dataclass}es (\code{LMParams},
        \code{OCPConfig}, \code{Scenario}).
  \item \textbf{Dynamics} --- the continuous model \code{flat\_moon\_6dof} and the
        integrators.
  \item \textbf{OCP build \& solve} --- \code{solve\_landing} and the analytic
        \code{cutoff\_freefall}.
  \item \textbf{Landing report} --- \code{print\_report}, \code{print\_cutoff\_report}.
  \item \textbf{Plots} --- states, controls, actuators, thrusters.
  \item \textbf{Trajectory plot} + \code{\_\_main\_\_} driver.
\end{enumerate}

\section*{1. Parameters}

Everything tunable lives in three dataclasses, so no magic numbers are buried in
the logic:
\begin{itemize}\itemsep2pt
  \item \code{LMParams} --- the vehicle: mass, inertia $J$, lunar gravity, DPS
        thrust limits $[T_{\min},T_{\max}]$, gimbal limit, the actuator models
        (thrust first-order lag $\tau_T$, gimbal second-order $\omega_n,\zeta$),
        and the 16-thruster RCS geometry. \code{rcs\_geometry()} returns the
        $6\times16$ allocation matrix $B$ mapping thruster forces to a body
        wrench; \code{B\_rcs} caches it (it is queried thousands of times while
        the graph is built).
  \item \code{OCPConfig} --- the problem: horizon \code{N}, step \code{dt},
        \code{integrator}, the path limits, the cost weights
        (\code{Qs,Qf,Rw,Rd}), and the cut-off altitude \code{h\_contact}.
  \item \code{Scenario} --- the initial state \code{x0} (position, velocity,
        attitude, rates) and the target \code{x\_target}.
\end{itemize}

\section*{2. State, dynamics and integrators}

The rigid body carries 12 states; five \emph{actuator} states are appended so the
optimiser sees the real, rate-limited hardware response:

\begingroup\small
\begin{center}
\begin{tabular}{@{}lll@{}}
\toprule
Index & Symbol & Meaning\\
\midrule
0--2   & $x_E,y_E,z_E$ & surface-frame position ($z_E\le0$ above ground)\\
3--5   & $u,v,w$       & body-frame velocity\\
6--8   & $\phi,\theta,\psi$ & Euler angles (3--2--1)\\
9--11  & $p,q,r$       & body rates\\
12     & $T$           & actual thrust (1st-order lag toward $T_c$)\\
13--14 & $\delta_p,\dot\delta_p$ & pitch gimbal (2nd-order toward $\delta_{p,c}$)\\
15--16 & $\delta_y,\dot\delta_y$ & yaw gimbal (2nd-order toward $\delta_{y,c}$)\\
\bottomrule
\end{tabular}
\end{center}
\endgroup

\code{flat\_moon\_6dof(x,u,lm)} returns $\dot{\bm x}$: navigation
($\dot{\bm r}=C_{E/B}\bm v$), body-frame translation (gravity + gimballed thrust
+ RCS force), Euler kinematics, Euler's rotational equations, and the actuator
ODEs. \code{rk4\_step} advances one step with classical RK4; \code{rk2\_step}
(Heun) is a cheaper alternative, and \code{INTEGRATORS} is the \code{\{'rk4','rk2'\}}
registry \code{solve\_landing} selects from.

\section*{3. Building and solving the OCP (\code{solve\_landing})}

The decision variables are the full state trajectory
$X\in\mathbb R^{17\times(N+1)}$ and the commands $U\in\mathbb R^{19\times N}$.
The constraints and objective are assembled on \code{Opti}:
\begin{itemize}\itemsep2pt
  \item \textbf{Initial condition} $X_{:,0}=\bm x_0$.
  \item \textbf{Dynamics (equality)} --- multiple-shooting defects
        $X_{:,k+1}=\Phi_{\mathrm{RK4}}(X_{:,k},U_{:,k})$ for every interval.
  \item \textbf{Path bounds} --- altitude, velocity, attitude, rate, thrust and
        gimbal limits on the states, and thrust/gimbal/RCS limits on the commands.
  \item \textbf{Cost} --- a stage term (state error $Q_s$ + command effort $R$),
        a \emph{control-rate} term $R_\Delta\lVert\Delta u_k\rVert^2$ that damps
        gimbal chatter, and a terminal term $Q_f$ pulling toward the contact
        condition ($z_E=-\code{h\_contact}$).
  \item \textbf{Warm start} --- states linearly interpolated $\bm x_0\!\to\!$
        target, thrust held at hover; a good guess cuts IPOPT iterations.
\end{itemize}
IPOPT runs with the \textbf{exact Hessian} (fewer, more informative iterations)
and \code{expand=True}, which rewrites the problem graph from MX to SX so the
derivatives evaluate about twice as fast. The \code{acceptable\_*} tolerances let
it stop at a good feasible point instead of grinding to the tight tol. The
function returns the solved \code{(Xs, Us)}.

\section*{4. Engine cut-off and touchdown (\code{cutoff\_freefall})}

The powered OCP deliberately ends a metre or so up, at the \emph{contact}
altitude. \code{cutoff\_freefall} then models the real DPS shutdown: it hard-cuts
the thrust ($T:=0$, gimbals zeroed, RCS off) and integrates the same 6-DOF
dynamics with fine \SI{0.02}{s} steps until $z_E=0$, interpolating the final
sub-step so touchdown sits exactly on the ground. For a near-rest contact the
touchdown speed is $v_{\mathrm{td}}=\sqrt{w_c^2+2gh_{\mathrm{contact}}}$. This
keeps the OCP well-posed near $z=0$ (no thrust gating, no hovering pathology).

\section*{5. Output}

\code{print\_report} summarises the powered terminal state (position, speed,
attitude vs. success thresholds); \code{print\_cutoff\_report} shows the cut-off
event and the analytic touchdown. Five plot helpers save PNGs: \code{plot\_states},
\code{plot\_controls} (DPS + net RCS wrench), \code{plot\_actuators}
(commanded vs. actual), \code{plot\_thrusters} (16 jets), and
\code{plot\_trajectory\_with\_axes} (3-D path with body-axis triads; the free-fall
segment is stitched on via \code{np.hstack([Xs, Xff])}).

\section*{Tuning \& performance knobs}

\begingroup\small
\begin{center}
\begin{tabular}{@{}lll@{}}
\toprule
Knob & Effect & Trade-off\\
\midrule
\code{N}, \code{dt}      & grid size / horizon & smaller $\Rightarrow$ faster solve, coarser near ground\\
\code{integrator}        & \code{rk4} / \code{rk2} & \code{rk2} cheaper/step but often more iterations\\
\code{expand}            & SX graph & $\sim\!2\times$ faster derivatives (kept on)\\
\code{h\_contact}        & cut-off altitude & lower $\Rightarrow$ softer touchdown $\sqrt{2gh}$\\
\code{Rd}[gimbal]        & move-suppression & higher $\Rightarrow$ smoother gimbals, stiffer NLP\\
\code{Qf}                & terminal accuracy & higher $\Rightarrow$ tighter landing, harder solve\\
\bottomrule
\end{tabular}
\end{center}
\endgroup

\noindent The dominant cost is IPOPT's KKT factorisation times the iteration
count, so \code{N} is the biggest speed lever; the trajectory becomes feasible
early and the remaining iterations polish optimality on a poorly-scaled problem.

\end{document}
